\documentclass[../../main.tex]{subfiles}% 注意这里的文件路径不能用 ./main.tex ,否则用latexmk编译子文件会报错
\graphicspath{{\subfix{./image/}}} % 指定图片目录,后续可以直接使用图片文件名
% 注意这里的文件路径不能用 ../../image/ ,否则用latexmk编译子文件会报错

% 例如:
% \begin{figure}[H]
% \centering
% \includegraphics[scale=0.3]{图.png}
% \caption{}
% \label{figure:图}
% \end{figure}
% 注意:上述\label{}一定要放在\caption{}之后,否则引用图片序号会只会显示??.

\begin{document}

\section{抽象曲面}

\begin{definition}
设$D\subseteq E^2$为二维欧氏空间中的开区域，若在$D$上给定正定二次微分形式
\begin{align}
\mathrm{I}=g_{\alpha\beta}\mathrm{d}u^\alpha\mathrm{d}u^\beta,\label{eq:6.0}
\end{align}
其中求和指标$\alpha,\beta=1,2$，$g_{\alpha\beta}=g_{\beta\alpha}$为$D$上的光滑函数，且矩阵$(g_{\alpha\beta})$在$D$内处处正定，则称$(D,\mathrm{I})$为\textbf{抽象曲面片}，$\mathrm{I}$称为该曲面片的\textbf{度量形式}。
\end{definition}

\begin{remark}
抽象曲面片的微分几何仅依赖度量形式$\mathrm{I}$，相关几何量包括Gauss曲率、曲线长度、测地曲率、测地线、测地坐标系、切向量平行移动等内蕴几何对象。
\end{remark}

\begin{definition}
设$M$为拓扑空间，若对任意$p\in M$，存在$p$的开邻域$U$、$\mathbb{R}^2$中的开区域$D$及同胚$\varphi:D\to U$，则称$M$为\textbf{二维拓扑流形}。
称$(U,\varphi)$为$M$的\textbf{坐标卡}；对$p\in U$，记
\begin{align}
u^\alpha(p)=(\varphi^{-1}(p))^\alpha,\quad \alpha=1,2,\label{eq:6.1}
\end{align}
则称$(u^1(p),u^2(p))$为点$p$的坐标，$(U;u^1,u^2)$为$M$的\textbf{坐标系}。
\end{definition}

\begin{definition}
设$M$为二维拓扑流形，若存在坐标卡集合$\{(U_i,\varphi_i):i\in I\}$满足
\begin{enumerate}[(1)]
\item $\{U_i:i\in I\}$构成$M$的开覆盖；
\item 对任意$i,j\in I$，若$U_i\cap U_j\neq\varnothing $，则坐标变换$\varphi_j^{-1}\circ\varphi_i$与$\varphi_i^{-1}\circ\varphi_j$均为$C^\infty$映射，
\end{enumerate}
则称$M$为\textbf{二维光滑流形}。
\end{definition}

\begin{definition}
设$M$为二维光滑流形，$p\in M$。若对包含$p$的任意容许坐标卡$(U_i,\varphi_i)$，存在数组$(v_i^1,v_i^2)$；且对任意两个包含$p$的容许坐标卡$(U_i,\varphi_i),(U_j,\varphi_j)$，满足
\begin{align}
v_j^\alpha=v_i^\beta\left.\frac{\partial u_j^\alpha}{\partial u_i^\beta}\right|_p,\quad \alpha=1,2,\label{eq:6.2}
\end{align}
则称该相容数组族为$M$在点$p$处的\textbf{切向量}。
\end{definition}

\begin{definition}
设$C_p^\infty$表示在$p$的某开邻域内有定义且各阶偏导数连续的函数集合。$M$在$p$处的切向量$\bm{v}$可等价定义为微分算子$\bm{v}:C_p^\infty\to\mathbb{R}$，满足
\begin{align}
\bm{v}(f)=v_i^\alpha\left.\frac{\partial(f\circ\varphi_i)}{\partial u_i^\alpha}\right|_{\varphi_i^{-1}(p)},\quad \forall f\in C_p^\infty,\label{eq:6.3}
\end{align}
其中$(U_i,\varphi_i)$为包含$p$的容许坐标卡。
\end{definition}

\begin{proposition}
切向量的微分算子定义\eqref{eq:6.3}与容许坐标卡的选取无关。
\end{proposition}

\begin{proof}

设$(U_j,\varphi_j)$为另一包含$p$的容许坐标卡，由链式求导法则：
\begin{align*}
v_j^\alpha\left.\frac{\partial(f\circ\varphi_j)}{\partial u_j^\alpha}\right|_{\varphi_j^{-1}(p)}
&=v_j^\alpha\left.\frac{\partial\left((f\circ\varphi_i)\circ(\varphi_i^{-1}\circ\varphi_j)\right)}{\partial u_j^\alpha}\right|_{\varphi_j^{-1}(p)}\\
&=v_j^\alpha\left.\frac{\partial(f\circ\varphi_i)}{\partial u_i^\beta}\right|_{\varphi_i^{-1}(p)}\left.\frac{\partial(\varphi_i^{-1}\circ\varphi_j)^\beta}{\partial u_j^\alpha}\right|_{\varphi_j^{-1}(p)}\\
&=v_i^\beta\left.\frac{\partial(f\circ\varphi_i)}{\partial u_i^\beta}\right|_{\varphi_i^{-1}(p)}.
\end{align*}
故算子取值不依赖坐标卡选取。

\end{proof}

\begin{proposition}
微分算子$\bm{v}:C_p^\infty\to\mathbb{R}$满足：
\begin{enumerate}[(1)]
\item 线性性：$\bm{v}(f+\lambda g)=\bm{v}(f)+\lambda\bm{v}(g)$；
\item 莱布尼茨法则：$\bm{v}(f\cdot g)=f(p)\bm{v}(g)+g(p)\bm{v}(f)$，
\end{enumerate}
其中$\forall f,g\in C_p^\infty,\lambda\in\mathbb{R}$。
\end{proposition}

\begin{proof}

由偏导数的线性性与乘积求导法则直接可得。

\end{proof}

\begin{definition}
对二维光滑流形$M$的坐标卡$(U_i,\varphi_i)$，定义切向量场$\frac{\partial}{\partial u_i^\alpha}:C^\infty(U_i)\to C^\infty(U_i)\ (\alpha=1,2)$：
\begin{align}
\frac{\partial}{\partial u_i^\alpha}(f)=\frac{\partial(f\circ\varphi_i)}{\partial u_i^\alpha},\quad \forall f\in C^\infty(U_i).\label{eq:6.4}
\end{align}
称$\left\{p;\left.\frac{\partial}{\partial u_i^1}\right|_p,\left.\frac{\partial}{\partial u_i^2}\right|_p\right\}$为$U_i$上的\textbf{自然标架场}。
\end{definition}

\begin{remark}
自然标架场在坐标域内处处线性无关，分别对应数组$(1,0),(0,1)$。
\end{remark}

\begin{definition}
设$M$为二维光滑流形，若对$M$的任意容许坐标卡$(U_i,\varphi_i)$，存在$U_i$上的正定二次微分形式
\begin{align}
g|_{U_i}=g_{\alpha\beta}^{(i)}\mathrm{d}u_i^\alpha\mathrm{d}u_i^\beta,\label{eq:6.5}
\end{align}
满足$g_{\alpha\beta}^{(i)}\in C^\infty(U_i),g_{\alpha\beta}^{(i)}=g_{\beta\alpha}^{(i)}$，且矩阵$(g_{\alpha\beta}^{(i)})$处处正定；且对交叠坐标卡$(U_i,\varphi_i),(U_j,\varphi_j)$，在$U_i\cap U_j$上有
\begin{align}
g_{\alpha\beta}^{(i)}=g_{\gamma\delta}^{(j)}\frac{\partial u_j^\gamma}{\partial u_i^\alpha}\frac{\partial u_j^\delta}{\partial u_i^\beta},\quad 1\leqslant\alpha,\beta\leqslant2,\label{eq:6.6}
\end{align}
则称$g$为$M$上的\textbf{度量形式}，其为2阶协变张量。
\end{definition}

\begin{definition}
若二维光滑流形$M$上指定了度量形式$g$，则称$(M,g)$为\textbf{二维黎曼流形}（抽象曲面）。
\end{definition}

\begin{proposition}
设$\bm{v}$为$M$在$p$处的切向量，在坐标卡$(U_i,\varphi_i)$下对应数组$(v_i^1,v_i^2)$，定义
\begin{align}
g|_p(\bm{v},\bm{v})=\left.g_{\alpha\beta}^{(i)}\right|_p v_i^\alpha v_i^\beta,\label{eq:6.7}
\end{align}
则该值与坐标卡选取无关，记$g|_p(\bm{v},\bm{v})=|\bm{v}|^2$。
\end{proposition}

\begin{proof}

设$(U_j,\varphi_j)$为另一坐标卡，$\bm{v}$对应$(v_j^1,v_j^2)$，由\eqref{eq:6.2},\eqref{eq:6.6}得：
\begin{align*}
\left.g_{\gamma\delta}^{(j)}\right|_p v_j^\gamma v_j^\delta
&=\left.g_{\gamma\delta}^{(j)}\right|_p \frac{\partial u_j^\gamma}{\partial u_i^\alpha}\frac{\partial u_j^\delta}{\partial u_i^\beta}v_i^\alpha v_i^\beta\\
&=\left.g_{\alpha\beta}^{(i)}\right|_p v_i^\alpha v_i^\beta.
\end{align*}
故取值与坐标卡无关。

\end{proof}

\begin{remark}
度量形式在坐标卡$(U_i,\varphi_i)$下的系数满足
\begin{align}
g_{\alpha\beta}^{(i)}=g|_{U_i}\left(\frac{\partial}{\partial u_i^\alpha},\frac{\partial}{\partial u_i^\beta}\right),\label{eq:6.8}
\end{align}
即$(g_{\alpha\beta}^{(i)})$为自然标架场的度量系数矩阵。
\end{remark}



\end{document}